\section{Introduction}
Cognitive behavioral therapy (CBT) is among the most heavily tested psychotherapies. A meta-analysis of 41 placebo-controlled trials in anxiety and related disorders (total $N=2{,}843$) found that CBT outperformed placebo on target symptoms with a moderate effect (Hedges' $g=0.56$) \citep{carpenter2018cbt}. The same literature also shows a ceiling: across 87 studies only about half of patients respond to CBT for anxiety disorders, and the figure depends on how response is defined \citep{loerinc2015response}; a larger meta-analysis of 100 articles reports a mean remission rate of 51.0\% \citep{springer2018remission}. A therapy can work on average and still leave open \emph{why} it works and \emph{for whom}.

The textbook account of CBT's cognitive component is a top-down loop. A situation triggers an automatic thought (``I am going to fail this talk''); the patient records it, names the distortion, weighs evidence for and against, and generates a more balanced thought; the feeling is expected to follow. This review first maps ten reasons given for why CBT fails (Section~\ref{sec:taxonomy}) and then grades the evidence for each. Its deepest case study is a narrow question about that loop: \emph{does the mechanism ask the patient to use a capacity that anxiety, and in some patients attention-deficit/hyperactivity disorder (ADHD), has already degraded?} If so, the cognitive component should depend on baseline executive capacity, should be harder to use when stressed than when calm, and should contribute less to outcome than behavioural elements that do not demand effortful reasoning.

The argument has a popular version. Judson Brewer, a clinician-researcher, argues that anxiety takes the prefrontal cortex ``off-line,'' so that willpower and reasoning are unavailable at the moments when they are needed \citep{brewer2016habit,brewer2021unwinding}, and he has applied this directly to CBT \citep{brewer2026cbt}. Our aim is to find out how much of this survives contact with the primary literature. Our answer, in brief, is that a narrow version survives and the strong version does not. The same standard is applied to the other reasons: dropout and non-response (Section~\ref{sec:nonresponse}), patient and context factors (Section~\ref{sec:patients}), therapist and delivery factors (Section~\ref{sec:delivery}), and the theory, measurement and evidence-quality critiques (Section~\ref{sec:mechanism}).

\paragraph{What we find} (1) Reappraisal does draw on working memory in laboratory studies, but evidence that acute stress abolishes it is thin and contested. (2) The direct evidence that executive function moderates CBT outcome is sparse and mostly null for overall outcome, with the clearest signal at the level of the cognitive skill itself. (3) Component studies repeatedly find that adding cognitive restructuring to exposure or behavioural activation does not add measurable benefit, but these studies are underpowered and a minority find a benefit. (4) CBT's headline effects shrink with active controls and bias adjustment; whether they shrink over time is disputed. (5) In adult ADHD, CBT is modestly effective, mostly as an adjunct to medication, and the plausible active ingredient is behavioural compensation; we found no trial testing ADHD as a moderator of CBT for anxiety in adults. (6) Beyond that argument, dropout (about 16\% before and 26\% during CBT), non-response, relapse and under-recorded harm are real and quantified, but mostly shared with other psychological treatments; patient, context and delivery factors shift outcomes without isolating a CBT-specific failure; and the theory of change, the outcome measures and the trial base are the weakest links.

\paragraph{Contributions} (i) An explicit decomposition of cognitive restructuring into executive demands (Section~\ref{sec:demands}); (ii) an evidence-graded review of the stress-and-prefrontal premise, including a check of the ``off-line'' claim against its sources (Section~\ref{sec:stress}); (iii) a review of efficacy, bias and component evidence (Sections~\ref{sec:evidence}--\ref{sec:ingredients}); (iv) a review of alternatives and of CBT in adult ADHD (Sections~\ref{sec:alternatives}--\ref{sec:adhdcbt}); (v) the direct moderator evidence (Section~\ref{sec:moderators}); (vi) a graded map of ten reasons CBT fails (Section~\ref{sec:taxonomy}) and reviews of non-response, patient, delivery and mechanism factors (Sections~\ref{sec:nonresponse}--\ref{sec:mechanism}); (vii) preregistrable predictions (Section~\ref{sec:agenda}).
